% \begin{itemize}
%     \item \underline{Стек-алгоритмы}
% \end{itemize}

% Рассмотрим один из возможных механизмов разрешения коллизий в рассмотренных системах.

% Предположим, что все терминалы слушают канал. В канале могут быть 3 вида событий:
% \begin{itemize}
%     \item успех --- когда в слоте передавал только 1 терминал,
%     \item коллизия --- когда в слоте передавало $\geqslant 2$ терминалов,
%     \item пустой слот.
% \end{itemize}

% Пусть исходно пыталось сделать передачу 4 терминала, начнём строить дерево коллизий; Терминалы, пришедшие в систему после начала построения дерева, ждут окончания его построения и уже после этого пытаются передавать.

% Построение дерева коллизий:

% \begin{center}
% \begin{tikzpicture}[every node/.style={font=\small}]
%     % Левая часть
%     \node[align=left] (left_list) at (0,0) {
%         $\bullet$ Успех \\
%         $\bullet$ Отсутствие \\
%         \hspace{0.3cm} информации
%     };
    
%     % Стрелка
%     \node at (2,0) {\huge $\longrightarrow$};
    
%     % Правая часть
%     \node[align=left] (right_list) at (4.5,0) {
%         $\bullet$ Успех \\
%         $\bullet$ Коллизия \\
%         $\bullet$ Пустой слот.
%     };
% \end{tikzpicture}
% \end{center}


% В момент изначального конфликта создается корневая вершина дерева, каждый из терминалов, который передавал в окне с корневой коллизией, с некоторой вероятностью выбирает передавать в следующем слоте или нет.

% Если терминал выбирает передавать, то он перемещается в левый потомок дерева, иначе --- в правый. Если конфликт происходит снова, то продолжаем строить дерево по тем же правилам. Если левая вершина оказалась концевой, то в следующем слоте начинает строиться правый потомок. Альтернативно, при наличии коллизии, сначала передаются обе ветви текущего уровня дерева, и только после этого, при наличии повторного конфликта, происходит дальнейшее ветвление. Однако среднее время разрешения изначальной коллизии в обоих вариантах одинаково.

% \begin{center}
% \begin{tikzpicture}[>=stealth, thick, 
%     slot_node/.style={draw, minimum width=1.2cm, inner sep=2pt, font=\small},
%     empty_node/.style={draw, dashed, minimum width=1.2cm, minimum height=0.5cm},
%     every node/.style={font=\small}
% ]

% % Номера этапов и узлы дерева
% \node at (0,0) {1.}; \node[slot_node] (s1) at (1.5,0) {1 2 3 4};
% \node at (0,-1.5) {2.}; \node[slot_node] (s2) at (1.5,-1.5) {1 2 4};
% \node at (0,-3) {3.}; \node[empty_node] (s3) at (0.5,-3) {};
% \node at (0,-4.5) {4.}; \node[slot_node] (s4) at (1.5,-4.5) {1 2 4};
% \node at (0,-6) {5.}; \node[slot_node] (s5) at (0.5,-6) {1};
% \node at (0,-7.5) {6.}; \node[slot_node] (s6) at (1.5,-7.5) {2 4};
% \node at (0,-9) {7.}; \node[slot_node] (s7) at (0.5,-9) {2};
% \node at (0,-10.5) {8.}; \node[slot_node] (s8) at (1.5,-10.5) {4};
% \node at (2.5,-11.5) {9.}; \node[slot_node] (s9) at (3.5,-11.5) {3};

% % Ребра (переходы)
% \draw[->] (s1.south) -- (s2.north);
% \draw[->] (s1.south) .. controls (4,-2) and (4,-10) .. (s9.north);

% \draw[->] (s2.south) -- (s3.north);
% \draw[->] (s2.south) -- (s4.north);

% \draw[->] (s4.south) -- (s5.north);
% \draw[->] (s4.south) -- (s6.north);

% \draw[->] (s6.south) -- (s7.north);
% \draw[->] (s6.south) -- (s8.north);

% % Текстовые пояснения
% \node[align=left, font=\small, text width=5cm] at (-4.5, -3) {так как слот пустой \\ и до этого была коллизия, то и \\ далее будет коллизия \\ $\Downarrow$ \\ можно сразу ветвиться};

% \node[align=center, font=\small, text width=6cm] at (1, -12.5) {закончил строиться левый потомок изначальной коллизии.};

% \node[align=left, font=\small, text width=7cm] at (7.5, -13) {начал (и сразу закончил) строиться правый потомок изначальной коллизии, можно не считать за коллизию, а вообще будем разрешать коллизию (если есть) правого потомка в следующем цикле.};

% \end{tikzpicture}
% \end{center}

% \begin{center}
% \begin{tikzpicture}[>=stealth, thick]
%     % Вспомогательная команда для отрисовки блоков пакетов
%     \newcommand{\packetbox}[3]{
%         \node[draw, anchor=south, inner sep=2pt, font=\small, minimum width=0.4cm] at (##1, 0) {\shortstack{##2}};
%     }

%     % Верхняя временная шкала (без оптимизации)
%     \draw[->, thick] (0,0) -- (10,0);
%     \foreach \x in {1,2,3,4,5,6,7,8,9} {
%         \draw (\x, 0.1) -- (\x, -0.1) node[below] {\x};
%     }
    
%     \packetbox{1}{1\\2\\3\\4}
%     \packetbox{2}{1\\2\\4}
%     \node at (3, 0.2) {+}; % Пустой слот (idle)
%     \packetbox{4}{1\\2\\4}
%     \packetbox{5}{1}
%     \packetbox{6}{2\\4}
%     \packetbox{7}{2}
%     \packetbox{8}{4}
%     \packetbox{9}{3}

%     % Стрелка перехода к оптимизации
%     \draw[double, ->, thick] (5, -1.2) -- (6, -2.2);

%     % Нижняя временная шкала (с оптимизацией пропуска повторной коллизии)
%     \begin{scope}[yshift=-4cm]
%         \draw[->, thick] (0,0) -- (8,0);
%         \foreach \x in {1,2,3,4,5,6,7} {
%             \draw (\x, 0.1) -- (\x, -0.1) node[below] {\x};
%         }
        
%         \packetbox{1}{1\\2\\3\\4}
%         \packetbox{2}{1\\2\\4}
%         \node at (3, 0.2) {+};
%         \packetbox{4}{1}
%         \packetbox{5}{2\\4}
%         \packetbox{6}{2}
%         \packetbox{7}{4}
%     \end{scope}
% \end{tikzpicture}
% \end{center}


% Однако, данный метод хоть и имеет очень высокую эффективность, но не используется из-за требования на канал обратной связи.


% Б. С. Цыбаков показал, что нормированная пропускная способность в системе синхронного случайного множественного доступа $\leqslant 0,587$, что грустно. Вместо того чтобы сокращать вероятности пустых и коллизионных слотов, давайте сокращать длительности простаивания канала и коллизий.

% $\Rightarrow$ 
% \begin{enumerate}
%     \item Прослушивание несущей.
%     \item Доступ по расписанию.
% \end{enumerate}

% $\rightarrow$ Если слот пустой, это становится ясно почти сразу: предполагаем, что все узлы находятся в зоне радиовидимости друг друга, время распространения сигнала между любыми двумя узлами $\ll$ длительности передачи пакета. Таким образом, прослушав канал в начале слота и не обнаружив передачу, узлы считают слот пустым. Так как до окончания слота ничего больше не произойдёт, то можно укоротить слот.

% Отличие между методами, которые будем рассматривать, состоит в действии, которое узлы должны предпринять для совершения попытки передачи пакета. Так как методы множественного доступа с прослушиванием предполагают, что узлы попеременно ведут приём/передачу, то введём предположение относительно их работы:

% время распространения сигнала между двумя самыми удалёнными станциями: $turn\_around\_time = \frac{l}{c}$, тогда длина слота $a > 2 \cdot \frac{l}{c}$, так как иначе у (1) будет сдвиг в логике протокола: для слота длины $a$ должна быть синхронность на MAC-уровне, то есть у двух станций может быть коллизия только если они выбрали один протокольный момент передачи.

% \begin{figure}[h!]
%     \centering
%     \begin{tikzpicture}[>=stealth, xscale=1.5, yscale=1.2]
%         % Оси времени
%         \draw[->] (0,1.5) -- (4,1.5) node[right] {$t_1$};
%         \draw[->] (0,0) -- (4,0) node[right] {$t_2$};
        
%         % Станции
%         \node[circle, draw, inner sep=2pt] at (-0.3, 1.5) {1};
%         \node[circle, draw, inner sep=2pt] at (-0.3, 0) {2};
        
%         % Фреймы/пакеты
%         \draw[thick] (0.2, 1.5) rectangle (1.2, 1.8) node[midway, font=\tiny] {frame 1};
%         \draw[thick] (0.2, 0) rectangle (1.2, -0.3) node[midway, font=\tiny] {frame 2};
        
%         % Линии распространения
%         \draw[dashed, ->] (0.2, 1.5) -- (1.0, 0);
%         \draw[dashed, ->] (0.2, 0) -- (1.0, 1.5);
        
%         % Обозначения моментов
%         \draw[->] (1.0, 1.9) -- (1.0, 1.6);
%         \draw (1.0, -0.1) -- (1.0, 0.1);
        
%     \end{tikzpicture}
% \end{figure}

% Кроме этого, слот включает в себя время прослушивания, которое зависит от формы сигнала: $a > 2 \cdot \frac{l}{c} + hardware\_delay + channel\_listen$.

% \textcolor{red}{Почему производительность ALOHA не зависит от $a$?}